\chapter{71}



«La voilà, Pasta mit Gemüse-Nusssauce und Schweizer Bergkäse.»

Louis balancierte die beiden Teller elegant nach draussen und stellte
sie einander gegenüber auf den Tisch. Mazi mischte gerade die Sauce
unter den Salat. Er blickte hoch.

«Wow, klingt vielversprechend.»

Er wedelte mit der Hand über seinen Teller, um den Duft einzufangen.

«Mhm, exzellent.»

Louis setzte sich.

«Ich hoffe, du magst es. Das Rezept ist von meiner Mutter.»

Louis’ Plan ging auf. Mazi schien mit seiner Erklärung zufrieden.
Wie vereinbart, verschwand Mazi schliesslich in Manuelas Praxis. Über
das offene Fenster zur Wiese drangen rhythmische Klänge hinaus. Ziemlich
laut und unverkennbar afrikanisch. Gelegentlich sang Mazi mit. Einzelne
Töne klangen schief. Louis rührte Mazis hemmungsloses Mitträllern.
Lächelnd räumte er den Tisch leer.

Das Hochsteigen ins \emph{Schäferloft} fühlte sich gut an. Louis
freute sich so sehr auf den intimen Moment mit seiner \emph{Gibson},
dass seine Finger zu flattern begannen. Ganz fein. Er konnte es kaum
erwarten, seine Gitarre hinaus zu tragen und sich hinzusetzen, mitten hinein 
in den kühlen Norden der \emph{Sunnewaid}. Der Platz war ihm wirklich lieb  
geworden. 

Er lehnte sich an die Hauswand, wischte sich die feuchten Hände an
der Hose ab und begann die Gitarre zu stimmen. Jede Saite liess er einzeln nachklingen. Unverzüglich lief noch einmal das Gespräch mit Valentina in ihm ab.
Seit sie von ihrer Befragung auf der \emph{Questur} erzählt hatte, stellte 
er sich immer wieder dieselbe Frage. Was machte er hier oben eigentlich? 
In Mailand war der Teufel los, während er auf der \emph{Sunnewaid} an der Hauswand stand. Und selten zwar, wie jetzt, hatte er das Gefühl, sich hier eine Auszeit zu leisten. Es war irre. Die Überleitung kam von alleine und plötzlich. Er sah seine Mutter vor sich. Viele, viele Jahre zurück. Sie hatte sich kleine Rituale geschaffen.
Es war während der ersten Monate in Mailand gewesen. Sie hatte vor ihrem Kaffee auf dem schmalen Balkon gesessen. Er musste etwa zwölf gewesen sein. Die kleine Wohnung war breit ausgefüllt mit «Ismael Serranos» Stimme, wie oft zu dieser Zeit, wenn er von der Schule nach Hause kam. Die Worte waren wieder da. 

«Was machst du gerade?» 

Inhaltlich war ihm einzig diese Frage geblieben. Die melancholische Melodie hatte so gar nicht zu Mamans
selbstzufriedenem Gesicht gepasst. Warum spanische Musik, hatte
er sich gefragt. Wie konnte sie nur. Nach allem, was mit Vater geschehen
war. Louis’ grenzenloser Stolz auf seine spanische Herkunft hatte sich
schon längst ins Gegenteil verkehrt. Mit der Zeit hatte er sich nur noch
genervt, wenn das Lied lief. Und eigentlich war ihm egal gewesen, was
sich seine Mutter anhörte. 

Louis legte die Gitarre neben sich ab. Wovon handelte eigentlich der Liedtext genau?
Er presste die Augen zusammen, zog das Handy aus der Hosentasche und
suchte nach dem Song. Er fand \emph{«Qué Andarás Haciendo»} als Live-Version. Beim Versuch, die Kopfhörer einzustecken, fielen sie ihm
aus der Hand. Es fühlte sich verwirrend an. Ob er sich wirklich in diese vergangene Zeit zurückdenken wollte?  
In die Geschichte seiner Mutter? Wozu eigentlich und warum jetzt?
Er hörte sich das Stück ein zweites, ein drittes Mal an und glaubte nach
und nach zu verstehen. Der Inhalt musste ihre Sehnsucht nach Leben widerspiegelt haben. 
«Vielleicht probierst du ein neues Kleid an. Planst eine Flucht, um das Meer zu sehen.»

\qrblock{Ismael Serrano \emph{«Qué Andarás Haciendo»}}{media/QR-Codes/041_Qué_Andarás_Haciendo(Live)_-_Include_speech_by_Ismael_Serrano_Ismael_Serrano_Spotify.png}

Louis fiel ein, dass sie vermehrt ausgegangen war, Kurse besucht hatte, deren Inhalte ihm nicht mehr
einfielen und von neuen Freundinnen gesprochen hatte. Ihr Verhältnis zueinander war zwar über die Jahre vertrauensvoll und
stabil geblieben. Aber war es nicht immer nur um ihn gegangen?
Wie wenig er noch heute von ihr wusste. Von ihren ganz persönlichen Dingen. 
Louis beendete das Lied. «Was tust du gerade? Eine Flucht planen, um das Meer zu sehen» - wie gut ihm die Vorstellung gefiel. 
Vielleicht wäre es besser gewesen, er hätte sich ganz woandershin abgesetzt. Viel weiter weg, in ein fernes Land, oder - 
plötzlich klingelte in seiner Nähe ein Handy. In sich versunken hatte er nicht wahrgenommen, dass Mazi 
unterdessen auf der Bank neben dem Eingang sass und in einem Buch las. Der zog nun das Gerät aus seiner Gesässtasche, 
nahm den Anruf entgegen und schoss unvermittelt hoch. 
